\section{Related Work}
\label{sec:related}

Structural Entropy is part of a broader family of graph-complexity and graph
clustering methods, but it differs from most classical objectives in treating
hierarchical organization as the primary object to be decoded.
We summarize the closest methodological neighbors below and use them as
reference points for the taxonomy and benchmark sections.

\paragraph{Modularity-based community detection.}
Modularity compares observed within-community edge density against a
configuration-model null graph~\cite{newman2004findingmodularity}.
Louvain and Leiden-style optimization methods make this objective highly
scalable, with Leiden additionally improving partition connectivity and
refinement stability~\cite{traag2019leiden}.
These methods are strong default baselines for topology-driven clustering,
but their flat partitions and known resolution effects make them less suited
to explicitly nested abstractions.
SE instead evaluates the coding cost induced by an encoding tree, so the
partition is interpreted as an information hierarchy rather than only a
density surplus.

\paragraph{Flow-based compression.}
The map equation and Infomap encode random-walk trajectories using module
codebooks~\cite{rosvall2008maps}.
They are conceptually close to SE because both explain graph organization
through compression, but the objects being compressed are different:
Infomap compresses flows, while SE compresses structural uncertainty under
volumes, cuts, and tree refinement.
This distinction matters in applications where the observed graph represents
an information system rather than an explicit navigation process.

\paragraph{Spectral partitioning.}
Spectral clustering and normalized cuts use eigenvectors of graph Laplacians
to expose low-conductance partitions~\cite{shi2000normalized,ng2002spectral}.
These methods provide strong relaxations for balanced graph cuts and remain
important for initialization and diagnostics.
However, spectral objectives are usually flat and externally parameterized by
the desired number of clusters.
SE can use spectral information, as in the VNE connection, but its target is a
minimum-description hierarchy rather than an eigenspace partition alone.

\paragraph{Deep graph clustering and pooling.}
Neural graph clustering methods such as DiffPool, MinCutPool, and DMoN learn
soft assignments inside GNN architectures~\cite{ying2018diffpool,bianchi2020mincutpool,tsitsulin2020dmon}.
They are effective when node attributes and task labels provide additional
semantic signal, but their objectives are often architectural regularizers
rather than standalone measures of structural organization.
Recent SE-based neural methods, including LSENet and DeSE, bridge this gap by
making structural entropy differentiable~\cite{sun2024lsenet,zhang2025dese}.
The journal contribution of this survey is to place these neural methods in
the same conceptual map as the original combinatorial SE algorithms and their
domain-specific extensions.

